\chapter{설계 결정과 형식화 순서}
\label{chap:plan}

\section{설계 결정}

\begin{description}
  \item[D1. 새 꼭짓점이 필요할 때] 보조정리~\ref{lem:split2}에서 \(\mathrm{Option}\,V\)를 쓰기로 확정했다.
    원래 꼭짓점 \(w\)는 \(\mathrm{some}\,w\), 새 꼭짓점은 \(\mathrm{none}\)이다.
    \texttt{SplitVertex.lean}은 유한하고 같음을 판정할 수 있는 \(V\)에서, 같은 끝점의 두 다트에서
    회전 체계를 구성한다. 서로 다른 두 다트가 같은 면에 있으면 오일러 지표 증가 정리를 적용한다. 원래 다트와 새 그래프 다트의 전단사로 회전과
    면 순열을 옮기며, 고립점과 성분 수를 직접 센다. \texttt{FaceBoundary.lean}이 이를 써서
    보조정리~\ref{lem:J1}의 면 경계 단사성을 증명한다. \texttt{FaceCycle.lean}은 면 순환을
    시작점과 정방향이 같은 \texttt{CyclicSequence}로 바꾼다(보조정리~\ref{lem:face_cycle}).
  \item[D2. 3-연결성의 대상] 고립점을 뺀 지지 집합 기준 (정의~\ref{def:conn3}). 수축이 고립점을 만들고,
    모든 그래프가 고정된 꼭짓점 타입 \(V\) 위에 있기 때문이다.
  \item[D3. 사이클 표현] 구간 논증에는 \(\mathbb{Z}/m \to V\) 수열이 편하다 (정의~\ref{def:cycseq}). 경로
    체계(정의~\ref{def:paths})를 만들 때만 Mathlib의 \texttt{Walk}로 바꾼다.
  \item[D4. 2-합 붙이기] 직접 증명했다 (보조정리~\ref{lem:sum2}, \texttt{GlueEdge.lean}).
    공통 변의 두 사본을 유지하는 중간 상태는 일반 다트 순열로 다루고, 중복 변을 지운 뒤
    \texttt{SimpleGraph}의 회전 체계로 옮긴다.
\end{description}

\section{형식화 순서}

아래의 모든 묶음을 형식화했다. 뒤의 묶음은 앞에서 만든 정의와 보조정리를 사용한다.
\begin{enumerate}
  \item \textbf{기본 도구와 경로}: 연결성, 변 수축, 순환 수열, 세분 경로 체계,
        작은 그래프와 꼭짓점 집합 넓히기, 세분과 경로 체계의 동치,
        순환 단어 분류(보조정리~\ref{lem:cyclic}).
  \item \textbf{붙이기와 세분}: 서로소·한 점·공통 변 붙이기(보조정리~\ref{lem:sum0},
        \ref{lem:sum1}, \ref{lem:sum2}), 공통 이웃과 교대 배치에서 금지 세분 구성
        (보조정리~\ref{lem:TK5_of_common}, \ref{lem:TK33_of_alt}),
        수축에서 세분 들어 올리기(보조정리~\ref{lem:lift}).
  \item \textbf{연결성과 축약}: 수축 가능한 변의 존재(정리~\ref{thm:contractible})와
        최소 반례의 3-연결성(정리~\ref{thm:reduce}).
  \item \textbf{면 수술}: 꼭짓점 삭제, 쪼개기, 면 경계 사이클과 이웃 순서
        (정의~\ref{def:delVertex}, 보조정리~\ref{lem:split2}, \ref{lem:J1},
        \ref{lem:face_cycle}, \ref{lem:J2}), 공통 이웃 없는 수축 되돌리기와 같은 면 변 넣기
        (보조정리~\ref{lem:uncontract}, \ref{lem:add_edge}), 이를 합친 꼭짓점 나누기
        (보조정리~\ref{lem:S}, \texttt{VertexSplit.lean}).
  \item \textbf{마무리}: 순환 단어를 실제 회전의 구간 조건과 사이클 위치로 옮기는 연결
        (보조정리~\ref{lem:neighbor_word}, \texttt{NeighborWord.lean}),
        변 수의 강한 귀납으로 어려운 방향과 필요충분조건
        (정리~\ref{thm:hard}, \ref{thm:kuratowski}, \texttt{Hard.lean}).
\end{enumerate}

\section{검사 실행}

\begin{verbatim}
lake build                               # Library and axiom audit
scripts/blueprint.sh                     # HTML/PDF, names, finite checks
python3 -I scripts/check_definition.py   # Planarity definition checks
\end{verbatim}
작업 규칙(공리, \texttt{\textbackslash planned}에서 \texttt{\textbackslash lean}으로 옮기기, 작성자와
검토자 분리)은 저장소의
\texttt{AGENTS.md}에 있다.
